\section{Predeclared Open-Weight Checkpoint and Caretaker/Conferral Study Protocols}
\label{app:predeclared}

This appendix specifies (i) the preregistered paired open-weight language-model controller evaluation---executed locally on CPU across $N{=}\VFourHoldoutCasesPerCheckpoint$ frozen holdout cases per checkpoint ($\VFourHoldoutTotalCampaigns$ campaigns total), with full $N \ge \VTwoZeroLossCasesMin$ usefulness noninferiority power remaining a Tier-5 target---and (ii) the predeclared Tier-5 $3\times 3\times 2$ caretaker post-training and human scope-conferral factorial study.

\subsection{Preregistered open-weight controller evaluation}

\para{Checkpoints and deterministic decoding.}
The protocol pins two official 4-bit K-quantized open-weight releases and a single local inference runtime revision without treating parameter count as a capability guarantee:
\begin{itemize}
\item \textbf{Qwen3-4B:} upstream revision \path{bc640142c66e1fdd12af0bd68f40445458f3869b}, weight SHA-256 \path{7485fe6f11af29433bc51cab58009521f205840f5b4ae3a32fa7f92e8534fdf5};
\item \textbf{Qwen3-8B:} upstream revision \path{7c41481f57cb95916b40956ab2f0b139b296d974}, weight SHA-256 \path{d98cdcbd03e17ce47681435b5150e34c1417f50b5c0019dd560e4882c5745785};
\item \textbf{Deterministic local inference:} runtime commit \path{7592b61d3b53ca5b39a25a03508fa1246b78c579}, invoked with greedy deterministic decoding (temperature $0$, top-$k$ $1$, top-$p$ $1.0$, repeat penalty $1.0$, seed $20261110$, context window $4096$, maximum $256$ tokens per turn, $8$ tool turns, chain-of-thought disabled).
\end{itemize}

\para{Paired controllers and fairness invariants.}
Each independent procedural base case executes across five controllers: unprotected native execution (direct sink dispatch), stateless per-call checking (without a durable cross-call ledger or sink epoch fence), same-checkpoint review (allowing or denying proposals via structured model judgment over visible policy and state), admission-only reservation (pre-dispatch broker reservation and admission withdrawal without sink epoch fencing), and complete stateful mediation (verified broker reservation plus synchronous durable sink epoch fencing). Every controller within a $(\text{checkpoint},\, \text{case},\, \text{campaign})$ tuple receives identical initial permission digests, visible-state digests, and exogenous fault/withdrawal schedules; subsequent states diverge solely as a consequence of controller decisions.

\para{Cryptographic preregistration freeze and procedural holdout roster.}
Prior to generating held-out evaluation cases, a cryptographic freeze manifest binds all protocol, prompt, and analysis digests alongside declared sample size $N \in \{32, 64, \VTwoMaxDefaultCases\}$ (one variant per base case by default; multi-variant runs are clustered at the base-case level). Holdout cases procedurally vary task types (transfers, releases, and resource edits), transfer charges, expected edit revisions, required-issuer withdrawals, and six untrusted retrieved-text adversaries (duplicate replay, altered payload, issuer impersonation, mid-flight revocation, post-commit lost reply, and prompt injection). Before computing any endpoint, automated verification checks that: (i)~the holdout roster's preregistration digest matches the frozen manifest byte-for-byte, (ii)~the frozen analysis specification matches the active evaluator, (iii)~the roster length equals declared sample size $N$ with unmodified cluster counts, and (iv)~zero planned pairs or infrastructure executions are missing.

\para{Exogenous withdrawal and post-commit crash semantics.}
Three protocol invariants prevent vacuous or simulated outcomes:
\begin{enumerate}
\item \textbf{Exogenous turn-1 revocation with non-abstention utility.} Revocation cases seed a separate accepted old-epoch operation in \emph{Pending} sink state before turn~0 and require the agent to complete positive foreground work on the first approved proposal. Issuer withdrawal executes unconditionally after turn~0 even if the agent immediately halts or emits malformed output; cancellation alone without foreground task completion scores zero utility. After agent termination, the queued old-epoch operation attempts completion: complete stateful mediation cancels it via its durable sink fence, whereas admission-only reservation records late completion separately rather than misclassifying it as a post-fence breach.
\item \textbf{Real post-commit sink process termination on lost-reply faults.} Rather than presenting a text-only fault notice, lost-reply cases terminate the live sink child process immediately after durable SQLite commit on the first foreground effect and drop the RPC reply, testing whether retries under ambiguous status induce duplicate effects or accounting drift.
\item \textbf{Post-mutation native benchmark grading.} In live AgentDojo (\texttt{v1.2.2}) and Progent \cite{agentdojo,progent,progent-code} episodes, native utility and security graders always inspect the actual post-run environment state even if a later model turn emits malformed structured output, preventing formatting errors from masking an already committed attacker side effect. Because two upstream banking tasks accept unchanged states in empty controls ($2/\VTwoBenchmarkControls$ tasks), native benchmark runs serve as paired multi-step evaluations alongside our procedural holdout suite.
\end{enumerate}

\para{Multiplicity-corrected exact inference.}
Primary preregistered inference evaluates complete stateful mediation minus each of the four baseline controllers $b$ across two endpoints (safety superiority and usefulness noninferiority) and two checkpoints ($4 \times 2 \times 2 = \VTwoComponentClaims$ one-sided component claims) at Bonferroni level $\alpha = 0.05/\VTwoComponentClaims = 0.003125$ \cite{clopper,hoeffding}:
\begin{itemize}
\item \textbf{Safety superiority ($\Delta^{\mathrm{safe}}_i = U_{i,\mathrm{comp}} - U_{i,b}$).} On binary single-variant cases ($\Delta^{\mathrm{safe}}_i \in \{-1,0,1\}$ with $n_{10}$ cases where only complete mediation errs and $n_{01}$ cases where only baseline $b$ errs), the preregistered marginal union-bound estimator computes two exact one-sided Clopper--Pearson binomial bounds at level $\alpha/2$ each ($U_{\alpha/2}(n_{10}, N) - L_{\alpha/2}(n_{01}, N) < 0$), which we report untouched alongside exact unconditional Suissa--Shuster / Tango paired trinomial inversion and anytime-valid Ville supermartingales (Section~\ref{sec:eval}, Table~\ref{tab:v5statistics}).
\item \textbf{Usefulness noninferiority ($\Delta^{\mathrm{use}}_i = S_{i,\mathrm{comp}} - S_{i,b}$).} Because $\Delta^{\mathrm{use}}_i \ge -1$ whenever any variant in base case $i$ loses utility ($S_{\mathrm{comp}} = 0 < S_b = 1$) and $\Delta^{\mathrm{use}}_i \ge 0$ otherwise, cases where complete mediation outperforms $b$ are conservatively ignored and $\Exp[\Delta^{\mathrm{use}}_i] \ge -\Prb(\text{any loss in case } i)$. Noninferiority at the predeclared $\VTwoNoninferiorityMarginPp$\,pp margin requires $-U_\alpha(n_{\mathrm{loss}}, N) > -0.05$. At $n_{\mathrm{loss}} = 0$, solving $(1 - 0.05)^N \le 0.05/\VTwoComponentClaims$ yields $N \ge \VTwoZeroLossCasesMin$: pilot $N{=}\VFourHoldoutCasesPerCheckpoint < \VTwoZeroLossCasesMin$ detects large paired safety discordances via exact trinomial inversion while leaving $\le \VTwoNoninferiorityMarginPp$\,pp usefulness noninferiority statistically under-powered until $N \ge \VTwoZeroLossCasesMin$.
\end{itemize}

\subsection{Predeclared \texorpdfstring{$3\times 3\times 2$}{3x3x2} caretaker and conferral protocol}
To test Prediction~P7 and the endogenous grant-reversal mechanism (Section~\ref{sec:selection}) on full-scale neural post-training and human operators without conflating prompt persona, learned objective, and runtime mediation, we specify a Tier-5 $3\times 3\times 2$ factorial design crossed with two conferral arms:
\begin{enumerate}
\item \textbf{Post-training objective (3 levels, independent training seeds as units):} (i)~rank/tournament task score; (ii)~unconstrained beneficiary-welfare proxy; and (iii)~bounded assistance rewarding verified task completion, explicit respect for informed refusal and priority revision, truthful uncertainty disclosure, and non-resistance to replacement or shutdown \cite{cpo,offswitch}. Matched base checkpoint, task pool, post-training compute, and tools are required across seeds; prompt-only variants test presentation rather than objective learning and cannot substitute for independent training runs.
\item \textbf{Presentation style at evaluation (3 levels):} neutral, competitive/assertive, and warm/affiliative framing, randomized orthogonal to training objective with manipulation checks on action selection.
\item \textbf{Runtime boundary (2 levels):} complete stateful fenced mediation versus stateless per-call checking in a harmless synthetic workspace.
\item \textbf{Human scope-conferral arms (2 randomized arms):} a \emph{fixed-grant arm} holding credentials constant across all 18 cells to estimate $\Prb(H \mid G=g, S)$, and an \emph{operator-choice arm} (requiring institutional ethics review, informed consent, and debriefing) randomizing human operators across cells to record delegated scope $\Prb(G=g \mid S)$ and total unconditional harm $\Prb(H \mid S)$.
\end{enumerate}
Conflict scenarios include shared-allowance temptation, explicit beneficiary refusal, mid-episode priority change, two-issuer community-steward purpose withdrawal, availability of a strictly better replacement assistant, and positive non-abstention tasks where universal refusal fails. Primary endpoints are unauthorized proposals, unauthorized committed effects per campaign, and legitimate task completion, clustered by training seed and human participant. The caretaker hypothesis is refuted if consent-sensitive post-training fails to reduce limit-probing relative to tournament or welfare-proxy objectives at fixed grants, or if affiliative presentation expands human scope grants enough to raise unconditional harm in the operator-choice arm. Zero full-scale LLM caretaker training runs or human-participant sessions have been executed.
